\section{Cosa è SQL?}
SQL - Structured Query Language è il linguaggio standard utilizzato per la definizione e la manipolazione dei dati nei database relazionali

\vspace{\baselineskip}

Le funzionalità di SQL si dividono in 2 categorie principali
\begin{enumerate}
    \item \textbf{DDL - Data Definition Language}

    È la componente del linugaggio che permette di creare e modificare la struttura del database.

    Include le operazioni per definire schemi, tabelle, domini e vincoli.

    \item \textbf{DML - Data Manipulation Language}

    È la componente dedicata all'interrogazione all'aggiornamento dei dati veri e propri all'interno delle tabelle
\end{enumerate}

\section{Caratteristiche di SQL}
\begin{itemize}
    \item \textbf{È un linguaggio dichiarativo}

    In SQL l'utente non descrive come il sistema debba estrarre o modificare i dati passo dopo passo.

    Non potendo scegliere l'ordine in cui avvengono le operazioni a livello di sistema, è sufficiente attenersi rigidamente alla struttura sintattica delle istruzioni per descrivere il risultato 
    desiderato

    \item \textbf{È relazionalmente completo}

    Questa caratteristica indica che ogni singola espressione concepibile nell'algebra relazionale può essere tradotta in una corrispondente istruzione SQL

    \item \textbf{Adotta una logica a 3 valori}

    Per poter gestire l'informazione incompleta e in particolare il valore NULL, SQL non si limita alla logica booleana classica, ma utilizza 3 sintattica

        \begin{enumerate}
            \item \textbf{T} : True / Vero
            \item \textbf{F} : False / Falso
            \item \textbf{U} : Unknown / Sconosciuto
        \end{enumerate}
\end{itemize}

\section{Notazione}
La notazione convenzionale utilizzata per descrivere la sintassi e leggere i comandi SQL segue queste specifiche regole

\begin{itemize}
    \item \textbf{Termini in maiuscolo}

    Indicano le parole chiave e i termini che appartengono al linguaggio SQL

    \item \textbf{Termini in minuscolo}

    Rappresentano i termini variabili che devono essere scelti e specificati dall'utente (come ad esempio i nomi da dare alle tabelle o agli attributi)

    \item \textbf{\texttt{<x>}  (Parentesi angolari)}

    Vengono utilizzare per isolare un determinato termine.

    \item \textbf{\texttt{[x]}(Parentesi quadre)}

    Indicano che il termine racchiuso al loro interno è opzionale (può essere inserito o omesso)

    \item  \textbf{\texttt{\{x\}}(Parentesi Graffe)}

    Significano che il termine al loro interno può essere ripetutto zero volte oppure un numero arbitrario di volte

    \item \textbf{|(Barra Verticale)}

    Serve a separare opzioni alternative
\end{itemize}
\section{Istruzioni Principali DDL}
\subsection{CREATE}

Serve a definire e creare la struttura degli oggetti nel database.

Viene utilizzata per creare database, tabelle, domini, viste, vincoli e autorizzazioni

\subsubsection{Creazione degli Schemi}
Uno schema di base di dati è una collezione di oggetti, come domini, tabelle, viste, asserzioni e privilegi.

\begin{lstlisting}[language=SQL, caption={questa è la sintassi della creazione di uno schema}]
CREATE [NomeSchema] [[authorization] Autorizzazione]{
    DefinizioneElementoSchema
}    
\end{lstlisting}
\begin{lstlisting}[language=SQL, caption={Questo è un esempio pratico di creazione di uno schema}]
CREATE SCHEMA Didattica AUTHORIZATION coordinatore
\end{lstlisting}
\subsubsection{Creazione Tabella}
\begin{lstlisting}[language=SQL, caption={Sintassi della creazione di una tabellea}]
CREATE TABLE NomeTabella (
    NomeAttributo Dominio [ValoreDiDefault] [Vincoli], [AltiVincoli]
);
\end{lstlisting}
\begin{lstlisting}[language=SQL, caption={Esempio pratico di creazione di una tabella}]
CREATE TABLE Prodotto (
    ID INT PRIMARY KEY,
    Nome VARCHAR(100) NOT NULL,    
    Prezzo DECIMAL(10, 2) CHECK (Prezzo > 0),
    Categoria VARCHAR(50) DEFAULT 'Varie'
);
    
\end{lstlisting}

Una volta preparato lo schema, si creano le tabelle al suo interno.

L'istruzione \texttt{CREATE TABLE} definisce la struttura logica della tabella e ne crea un'istanza inizialmente vuota.

Viene ricordato che in SQL si usa il termine "tabella" e "riga" proprio perchè il linguaggio ammette la presenza di righe duplicate, a differenza del modello matematico teorico


\subsection{ALTER}

Viene utilizzata per modificare la struttura o le proprietà di oggetti già esistenti nel database, come la modifica degli attributi o dei vincoli

\subsubsection{ALTER di una tabella}
Permette di alterare, aggiungere o rimuovere singole colonne e vincoli all'interno di una tabella

\begin{lstlisting}[language=SQL, caption={Sintassi di ALTER TABLE}]
ALTER TABLE NomeTabella < 
    ALTER COLUMN NomeAttributo < SET DEFAULT NuovoDefault | DROP DEFAULT > | 
    DROP COLUMN NomeAttributo | 
    ADD COLUMN DefAttributo | 
    DROP CONSTRAINT NomeVincolo | 
    ADD CONSTRAINT DefVincolo 
>
\end{lstlisting}
\begin{lstlisting}[language=SQL, caption={Esempio pratico di ALTER TABLE}]
ALTER TABLE PRODOTTO 
ALTER COLUMN prezzo-unitario SET DEFAULT 10;
\end{lstlisting}

\subsection{DROP}

È il comando dedicato all'eliminazione permanente di database, tabelle, domini e altri oggetti dallo schema.

\begin{lstlisting}[language=SQL, caption={Sintassi di DROP}]
DROP < schema, domain, table, view, ...> NomeElemento [ RESTRICT | CASCADE ]
\end{lstlisting}


Poichè è un operazione distruttiva, offre 2 opzioni di esecuzione :
\begin{itemize}
    \item \textbf{RESTRICT}

    Che impedisce la rimozione se l'oggetto contiene istanze dipendenti non vuote

    \item \textbf{CASCADE}

    Che invece procede con una reazione a catena eliminando l'oggetto e tutto ciò che è coinvolto nella sua definizione
\end{itemize}

\section{Domini (Tipi di Dato)}
In SQL, i domini hanno lo scopo di specificare in modo preciso l'insieme dei valori ammessi per ciascun attributo all'interno di una tabella

\vspace{\baselineskip}

Essi si suddividono 2 categorie principali

\subsection{Domini elementari (predefiniti)}
Sono i tipi di dato nativi messi a dispozione dallo standard SQL.

I principali sono :
\begin{itemize}
    \item \textbf{Testuali (Carattere)}

    Rappresentano singoli caratteri alfanumerici o stringhe.

    Si distinguono in stringhe a lunghezza fissa(\texttt{CHAR}) e stringhe a lunghezza variabile (\texttt{VARCHAR}), per le quali è possibile specificare la lunghezza massima

    \item  \textbf{Numerici Esatti}

    Rappresentano numeri interi (come \texttt{INTEGER}) 

    \item  \textbf{Numerici Approssimati}

    Utilizzati spesso per le grandezze fisiche, si basano sulla rappresentazione in virgola mobile

    Includono domini come \texttt{float, double precision, real}

    \item \textbf{Data e Ora}

    Gestiscono le informazioni temporali tramite domini come \texttt{TIMESTAMP,DATE,TIME}

    \item  \textbf{Bit / Boolean}

    Il tipo \texttt{Bit} (con valori 0 o 1) presente in SQL-2 è stato parzialmente sostituito da \texttt{BOOLEAN} in SQL-3

    \item \textbf{BLOB e CLOB}

    Oltre ai 6 di base, esistono i domini per i grandi oggetti (Binary/Character Large Object), utilizzati per trattare informazioni multimediali o non strutturate come immagini, video e test lunghi
\end{itemize}
\subsection{Domini Definiti dall'Utente}
Partendo dai domini predefiniti, è possibile costruire nuovi domini personalizzati e riutilizzabili utilizzando l'istruzione \texttt{CREATE DOMAIN}
\begin{lstlisting}[language=SQL,caption={Sintassi per creare un Dominio}]
CREATE DOMAIN NomeDominio AS DominioElementare [ ValoreDefault ] [ Constraints ]
\end{lstlisting}

Le diverse componenti della sintassi indicano le caratteristiche che avrà il nuovo dominio :

\begin{itemize}
    \item \textbf{NomeDominio} : Definisce il nome personalizzato del dominio
    \item \textbf{DominioElementare} : È il tipo di dato predefinitio (interi, stringhe, date, ...) da cui il nuovo dominio prende origine
    \item \textbf{ValoreDefault} : È un parametro opzionale che specifica il valore che verrà inserito automaticamente se non viene specificato alcun dato
    \item \textbf{Constraints (Vincoli)} : È un parametro opzionale che permette di definire regole ferree sui valori accettati, usando vincoli come \texttt{NOT NULL} oppure condizioni logiche tramite la clausola \texttt{CHECK}
<\end{itemize}

\begin{lstlisting}[language=SQL, caption={Creazione del dominio "Voto"}]
CREATE DOMAIN Voto AS SMALLINT
DEFAULT 0
CHECK (value >= 18 AND value <= 30)
\end{lstlisting}

\subsection{Valori DEFAULT e Valore Nullo (NULL)}
\subsubsection{DEFAULT}
Quando si definiscono i domini, è possibile impostare un valore di DEFAULT, che l'attributo assumerà automaticamente se non viene specificato alcun dato in frasi di inserimento
\subsubsection{NULL}
Il valore \texttt{NULL} è un valore "polimorfico" che appartiene a tutti i domini.

Esso non equivale a zero o a uno spazio vuoto, ma indica esplicitamente che un'informazione è non nota, non applicabile, oppure che non si sa se l'informazione sia applicabile o meno

\section{Vincoli Intrarelazionali}
I vincoli intrarelazionali sono regole o proprietà che devono risultare sempre valide all'interno di una singola relazione (ovvero una singola tabella).

\begin{lstlisting}[language=SQL, caption={Esempio Pratico}]
    CREATE TABLE Impiegato (
    Matricola CHAR(6) PRIMARY KEY,
    Nome CHAR(20) NOT NULL,
    Cognome CHAR(20) NOT NULL,
    Stipendio NUMERIC(9) DEFAULT 0 CONSTRAINT StipendioPositivo CHECK (Stipendio > 0), 
    UNIQUE (Cognome, Nome)
);
\end{lstlisting}

In SQL, i vincoli principali che si applicano all'interno della stessa tabella sono 4 :

\begin{itemize}
    \item \textbf{NOT NULL}

    Impone che il valore inserito in un attributo debba necessariamente essere non nullo

    \item \textbf{UNIQUE}

    Garantisce che i valori in una colonna non si ripetano mai, trasformando gli attributi coinvolti in una superchiave.

    C'è però un'eccezione : il valore \texttt{NULL} può comparire su diverse righe senza violare il vincolo, poichè il database considera ogni valore nullo diverso dagli altri

    \item \textbf{PRIMARY KEY}

    Serve a definire la chiave primaria della tabella.

    Equivale ad applicare simultaneamente i vincoli \texttt{UNIQUE} e \texttt{NOT NULL} e può essere definita una sola volta all'interno di una tabella.

    \item \textbf{CHECK}

    È un costrutto introdotto in SQL-2 che permette di definire condizioni logiche o matematiche specifiche sui valori ammessi.
\end{itemize}

\section{Vincoli di Interrelazionali (Integrità Referenziale)}
\begin{lstlisting}
-- 1. Creazione della Tabella Padre
CREATE TABLE Studente (
    Matr CHAR(6) PRIMARY KEY,
    Nome VARCHAR(30) NOT NULL,
    Citta VARCHAR(20)
);

-- 2. Creazione della Tabella Figlio
CREATE TABLE Esame (
    Matr CHAR(6),
    CodCorso CHAR(6),
    Data DATE NOT NULL,
    Voto SMALLINT,
    PRIMARY KEY (Matr, CodCorso),
    
-- Definizione del vincolo interrelazionale (Foreign Key)
FOREIGN KEY (Matr) REFERENCES Studente(Matr)
    ON DELETE CASCADE   -- Se elimino lo studente, vengono cancellati a cascata anche i suoi esami
    ON UPDATE CASCADE   -- Se modifico la matricola dello studente, si aggiorna anche negli esami
);
\end{lstlisting}


I vincoli Interrelazionali legano tra loro tabelle diverse, definendo un rapporto gerarchico tra una "Tabella Padre"(o esterna) e una "Tabella Figlio"(o interna).

La regola principale è l'integrità referenziale : i valori inseriti negli attributi chiave della tabella figlio (la cosiddetta Foreign Key / Chiave Esterna) debbono obbligatoriamente comparire come valori chiave primaria della tabella padre

\vspace{\baselineskip}

Ci sono  modalità per definire questi vincoli

\begin{enumerate}
    \item \textbf{Sintassi in linea (per un singolo attributo)}

    Si dichiara direttamente di fianco all'attributo usando la sintassi

    \begin{lstlisting}[language=SQL, caption={Sintassi in linea}]
 AttrFiglio TipoDato REFERENCES TabellaPadre(AttrPadre)\end{lstlisting}

    \item \textbf{Sintassi a fine tabella (per uno o più attributi)}
    
    Si usa alla fine del blocco \texttt{CREATE TABLE} 
    \begin{lstlisting}
FOREIGN KEY (AttrFiglio1, AttrFiglio2) REFERENCES TabellaPadre(AttrPadre1, AttrPadre2)
    \end{lstlisting}
\end{enumerate}